\documentclass{article}

% NeurIPS 2026 main-track submission format.
% Do not add "final" or "preprint" for an anonymized submission.
\usepackage{neurips_2026}

\usepackage[utf8]{inputenc}
\usepackage[T1]{fontenc}
\usepackage{hyperref}
\usepackage{url}
\usepackage{booktabs}
\usepackage{amsfonts}
\usepackage{amsmath}
\usepackage{amssymb}
\usepackage{graphicx}
\usepackage{nicefrac}
\usepackage{microtype}
\usepackage{xcolor}

\title{Microglia-Inspired Dynamic Pruning for Reasoning Models}

\author{%
  Anonymous Author(s) \\
  Affiliation \\
  Address \\
  \texttt{email} \\
}

\begin{document}

\maketitle

\begin{abstract}
  Replace this paragraph with a one-paragraph summary of the problem, method,
  main empirical result, and scope of the contribution.
\end{abstract}

\section{Introduction}

State the problem, the proposed microglia-inspired dynamic pruning approach, and
the main claims supported by the experiments.

\section{Related work}

Summarize related work on transformer pruning, adaptive inference, sparse
attention, and biologically inspired neural computation.

\section{Method}

Describe the activation statistics, microglia agents, masking mechanism, and
training objective.

\section{Experiments}

Describe models, datasets, baselines, matched sparsity controls, evaluation
metrics, hardware, seeds, and statistical tests.

\section{Results}

Report accuracy, latency, sparsity, confidence intervals, ablations, and mask
behavior analyses.

\section{Limitations}

Discuss limits of the empirical setting, model scale, datasets, compute budget,
and assumptions behind the pruning mechanism.

\begin{ack}
Acknowledgments and funding disclosures are hidden during anonymous submission
and appear only in final mode.
\end{ack}

\section*{References}

{\small
Add references here or replace this section with BibTeX commands.
}

\appendix

\section{Technical appendix}

Add supplementary experimental details, proofs, tables, or figures here.

\newpage
\input{checklist.tex}

\end{document}
